\begin{solucion}
    Por simple inspección, proponemos la siguiente posibles raíces $r=\pm 1$ y $r=\pm 3$. Y evaluando cada una de estas en polinomio obtenemos:
    \begin{align*}
    f(1) &= 1^3 - 7(1)^2 + 3(1) + 3 = 1 - 7 + 3 + 3 = 0 \\
    f(-1) &= (-1)^3 - 7(-1)^2 + 3(-1) + 3 = -1 - 7 - 3 + 3 = -8 \\
    f(3) &= 3^3 - 7(3)^2 + 3(3) + 3 = 27 - 63 + 9 + 3 = -24 \\
    f(-3) &= (-3)^3 - 7(-3)^2 + 3(-3) + 3 = -27 - 63 - 9 + 3 = -96 \\
    \end{align*}
    Asi por el teorema de factorización única tenemos que:

    \begin{align*}
        t^3 - 7t^2 + 3t + 3 &= (t-1)f(t) \\
    \end{align*}

    Donde $f(t)$ es un polinomio de grado 2. Para encontrarlo, realizamos la división sintética:

    \begin{align*}
        \begin{array}{r|rrrr}
            1 & 1 & -7 & 3 & 3 \\
              &   & 1 & -6 & -3 \\
            \hline
              & 1 & -6 & -3 & 0 \\
        \end{array}
    \end{align*}
    Por lo tanto, el polinomio $f(t)$ es:
    \begin{align*}
        f(t) &= t^2 - 6t - 3
    \end{align*}
    Luego el polinomio por la definición de irreducibilidad el polinomio $t^3 - 7t^2 + 3t + 3$ es irreducible, puesto que:
    \begin{align*}
        t^3 - 7t^2 + 3t + 3 &= (t-1)(t^2 - 6t - 3) \\
        &= (t-1)(t^2 - 6t - 3) \\
    \end{align*}
    Es decir, el polinomio $t^3 - 7t^2 + 3t + 3$ es producto de dos polinomios de grado menor. Lo cual es valido pues $\mathbb Q$ es un cuerpo.
\end{solucion}